% Generated draft: requires mathematical and source review.
\documentclass[11pt]{article}
\usepackage[margin=25mm]{geometry}
\usepackage{fontspec}
\setmainfont{Latin Modern Roman}
\usepackage{amsmath,amssymb,mathtools,unicode-math}
\setmathfont{Latin Modern Math}
\usepackage{enumitem,needspace,xurl,hyperref}
\setlength{\parindent}{0pt}
\setlength{\parskip}{5pt}
\setlength{\emergencystretch}{3em}
\newcommand{\N}{\mathbb{N}}
\newcommand{\Z}{\mathbb{Z}}
\newcommand{\Q}{\mathbb{Q}}
\newcommand{\R}{\mathbb{R}}
\newcommand{\C}{\mathbb{C}}
\newcommand{\F}{\mathbb{F}}
\newcommand{\A}{\mathbb{A}}
\newcommand{\PP}{\mathbb{P}}
\newcommand{\E}{\mathbb{E}}
\newcommand{\eps}{\varepsilon}
\newcommand{\dd}{\,\mathrm{d}}
\begin{document}
\section*{UnsolvedMath ID: 701309}
\textbf{Source problem code:} catalogue entry
\subsection{Definitions and mathematical statement}
Let \(G\) be an arbitrary group with identity \(1\). Use the commutator convention \([x,y]=x^{-1}y^{-1}xy\), and define the iterated commutators by
\[
[a,{}_1 g]=[a,g],\qquad [a,{}_{n+1}g]=[[a,{}_n g],g]\quad(n\geq 1).
\]
An element \(a\in G\) is a \emph{right Engel element} if
\[
\forall g\in G\;\exists n=n(a,g)\geq 1:\quad [a,{}_n g]=1.
\]
It is a \emph{bounded right Engel element} if
\[
\exists n=n(a)\geq 1\;\forall g\in G:\quad [a,{}_n g]=1.
\]
Thus, in the bounded condition, the bound may depend on the element \(a\), but is uniform over \(g\in G\). Write \(R(G)\) and \(\overline{R}(G)\) for these two sets, respectively.

\begin{enumerate}
\item[(a)] Does there exist a group \(G\) such that \(R(G)\) is not a subgroup of \(G\)?
\item[(b)] Does there exist a group \(H\) such that \(\overline{R}(H)\) is not a subgroup of \(H\)?
\end{enumerate}
The groups in the two questions need not be the same. No finiteness or other restriction on the groups is stated. The supplied record does not explicitly define the Engel terminology; the definitions above state the standard interpretation used in this draft.
\subsection{Short English statement}
A. Abdollahi asks whether the right Engel elements of some group fail to form a subgroup, and separately whether the bounded right Engel elements of some group fail to form a subgroup. In the bounded condition, each element has its own bound, valid against every other group element. The registry labels the problem open; this is supplied source information, not an independently verified status assessment.
\subsection{Sources}
\begin{itemize}
\item Local UnsolvedMath registry, supplied complete source record: permanent ID \(701309\), exact source code \texttt{catalogue entry}, title \emph{Kourovka Notebook Problem 17.7}. The record attributes the problem to A. Abdollahi, with proposed year 2010, and labels it \texttt{open}.
\item Supplied source attribution: \emph{The Kourovka Notebook (Unsolved Problems in Group Theory)}, 17th issue (2010), Problem 17.7; also identified in the record as appearing in the 21st issue edition (2026), \texttt{arXiv:1401.0300v48}. The supplied collection description names E. I. Khukhro and V. D. Mazurov as editors of the collection covering Issues 1--21 (1965--2026).
\item Registry revision: \texttt{d371ee917ac9647fbb06cd6fc529deec6f3dda82}; record updated \texttt{2026-10-16T14:00:00Z}. This is a registry timestamp, not a supplied literature-verification date. The cited editions and problem status have not been independently checked for this draft.
\end{itemize}
\end{document}
